\slajdid{02-s033}
\begin{frame}[label=02-s033]
Mnogi algoritmi minimizacije su pogodni za mašinsku implementaciju ali nisu baš za ručnu sintezu.

Najpoznatiji algoritam za ručnu minimizaciju je minimizacija putem Karnoovih tabela, karti, koja je limitirana sa brojem ulaznih promenljivih za koje ima smisla raditi minimizaciju na taj način.

Do 4 ulazne promenljive jako jednostavno, radi se sve u 2D prostoru, sa 5 promenljivih se već prelazi u 3D prostor, sa 6 promenljivih u 4D prostor gde već postoji problem da se ispravno prati proces minimizacije, a za 7 već praktično postaje neupotrebljiv.
\end{frame}
